\textbf{а)} Рассмотрим $\frac{a}{b} \in Quot(Quot(K))$, где $a, b \in Quot(K)$. Т.к. $Quot(K)$ -- поле, то $\frac{a}{b} \in Quot(K)$. Тогда можно поставить в соответствие элемент $\frac{(\frac{a}{b})}{e}$ (где $e$ -- единичные элемент) из $Quot(Quot(K))$ и $\frac{a}{b}$ из $Quot(K)$. Получим изоморфизм.

\textbf{б)}

Воспользуемся универсальным свойством. Но для этого сначала его напомниим:

Пусть у нас есть вложение $f: K \to L$, $L$ -- поле. Также у нас есть вложение $i: K \to Quot(K)$. В этом случае существует единственное вложение $\phi: Quot(K) \to L$, которое продолжает $f$, то есть $\phi \circ i = f$.

Легко понять, как строится это вложение:

Воспользуемся тем, что $L$ поле. Тогда если $f(a) \in L \Rightarrow f(a)^{-1} \in L$, поэтому положим $\phi(\frac{a}{b}) = f(a)f(b)^{-1}$\\

Теперь вернемся к нашей задаче. Можно построить следующие вложения: $i_1: K \to L, i_2: L \to Quot(K), \alpha: K \to Quot(K), \beta: L \to Quot(L)$. Теперь несколько раз воспользуемся универсальным свойством, а именно можно построить инъекцию $\psi: Quot(K) \to Quot(L)$ (в качестве отображения $f$ берем $\beta \circ i_1$), а также $\phi: Quot(L) \to Quot(K)$. В итоге получим две инъекции: $\phi, \psi$ между $Quot(K)$ и $Quot(L)$. Значит, они изоморфны.

